% Chapter 3.7: Boundary Condition Enforcement: Neural Operators vs. Classical Solvers
% Part III of the book.
% STATUS: rewritten to book standard (Part III draft pass).

\chapter{Boundary Condition Enforcement: Neural Operators vs. Classical Solvers}
\label{ch:part3-boundary-condition-enforcement-neural-oper}
\label{sec:boundary_conditions}
% -----------------------------------------------------------------------------

A boundary condition is half of the physical statement, not a comment on it. In every application
of this part the trace of the solution on some surface carries either the physics or the
deliverable.
\index{boundary condition!enforcement regimes}
External aerodynamics is the plainest case: no-slip
\index{no-slip condition}
on the body is a Dirichlet condition on the velocity, and the reported quantities of
interest---drag and lift coefficients, surface pressure, wall shear stress (WSS)---are integrals
over the same surface (\chapref{ch:part3-external-aerodynamics-and-surface-flows}). In cardiac
electrophysiology the zero-flux condition through the endocardium and epicardium helps determine
whether a reentrant spiral wave persists, drifts, or annihilates
(\chapref{ch:part3-cardiac-electrophysiology-and-arrhythmia-d}). In a faulted aquifer the
load-bearing constraint sits inside the domain, on the impermeable caprock across which carbon
dioxide has no normal flux (\chapref{ch:part3-subsurface-multiphase-flow-and-carbon-stor}).
Global weather prediction supplies the exception that defines the rule: the horizontal domain is
the sphere, which has no boundary, while the vertical has a ground and a model lid, and what the
poles require is regularity rather than boundary data
(\chapref{ch:part3-global-weather-forecasting-and-climate-dyn}). Tomography has no boundary
condition at all but has the analogous object, the support of the attenuation field together with
the detector's sampling domain (\chapref{ch:part3-medical-imaging-and-tomographic-reconstruc}).

This chapter treats that structure in general, and it belongs between the application chapters
because the comparison the book has been building is sharpest at the boundary. A conforming
discretization cannot return a field that violates its boundary conditions, because the condition
is a property of the space being searched; no spectral, graph, or branch-and-trunk layer has that
property by construction. A trained operator reproduces boundary values when its labels carry them
and fails when they do not, and the resulting error is nearly invisible in an average norm while
decisive in a certified design loop. \chapref{ch:part3-the-multi-query-bottleneck-in-computationa}
established why the surrogate is worth building. Added here is an accounting of what is given up,
and of the three ways to buy it back: build the condition into the parameterization, add it to the
objective, or leave it to the data and measure the consequence.

Section~\ref{sec:ch37-classical} states the classical mechanisms and pins down what ``hard'' means.
Section~\ref{sec:ch37-operator} gives the two operator-side mechanisms and the
qualifications each carries, Section~\ref{sec:ch37-structural} collects the structural
routes---coordinate channels, lifting, parameterizations that satisfy the condition by
form---and their cost in regularity, geometry work, and guarantees, and Section~\ref{sec:ch37-ledger}
prices the mechanisms
against an elliptic stability estimate before separating what is proven from what is observed.

% -----------------------------------------------------------------------------
\section{Classical Numerical Paradigms (FEM, FVM, Spectral)}
\label{sec:ch37-classical}

The three workhorses of scientific computing enforce boundary conditions by one device, and the
device is not a penalty. In a conforming discretization the condition is imposed on the space from
which the numerical solution is drawn, so every member of that space---the exact discrete solution,
every Krylov or Newton iterate on the way to it, and every diagnostic read off them---satisfies it
by construction. Nothing is appended to an objective and no weight is tuned. The three variants
named in the section title are three realizations of that idea.

Consider the finite element method (FEM)
\index{finite element method (FEM)!essential boundary conditions}.
Trial and test spaces are chosen before any basis, quadrature, or solver appears; the Darcy problem
\eqref{eq:darcy} was put in that form as \eqref{eq:ch03-weakform}, with test space $H^1_0(\Omega)$,
the closure of $C_c^{\infty}(\Omega)$ in $H^1(\Omega)$ from
\chapref{ch:part1-function-spaces-duality-and-sobolev-regula}. For the Poisson model the statement
is
\begin{equation}
  \label{eq:ch37-fem-weak}
  \int_\Omega \nabla u \cdot \nabla v \, dx = \int_\Omega f v \, dx + \int_{\partial\Omega} g v \, ds
  \qquad \text{for all } v \in H^1(\Omega) \text{ with } v|_{\Gamma_D}=0,
\end{equation}
where $\Gamma_D$ is the part of $\partial\Omega$ carrying prescribed values and $g$ the prescribed
Neumann (traction) data on the complementary part $\Gamma_N$; the boundary integral is written over
all of $\partial\Omega$ and reduces to $\Gamma_N$, because the test function vanishes on $\Gamma_D$.
With $\Gamma_D=\partial\Omega$ the display collapses to the pure-Dirichlet form,
\eqref{eq:ch03-weakform} with $a\equiv1$; with $\Gamma_D=\varnothing$ it is the pure-Neumann
problem, which is
solvable only when $f$ and $g$ satisfy a compatibility condition and then only up to a constant.
Two facts about
\eqref{eq:ch37-fem-weak} are habitually collapsed into one. The Dirichlet condition carried by the
trial space is \emph{essential}: it is imposed, and it holds as an identity of the linear system,
which is why Poincar\'e's inequality \eqref{eq:ch02-poincare}, in its variant for functions
vanishing on a part of the boundary, makes the energy gradient a
norm \index{boundary condition!essential and natural}. The Neumann condition is \emph{natural}: it
enters as data in the right-hand functional and is satisfied only in the integrated sense against
test functions, so the prescribed flux is recovered from the discrete solution rather than built
into it~\cite{evans2010partial,quarteroni2008numerical}. The distinction matters because it
identifies which half of a boundary condition a change of representation preserves for free, and it
bounds the word ``exact'': a curved $\partial\Omega$ is replaced by straight or isoparametric
faces, the condition is imposed on that approximation, and the discrepancy is the standard geometric
error of the method \index{finite element method (FEM)!geometric approximation error}, removable
only by refinement and only where the boundary is smooth enough for the interpolation to converge.

The finite volume method (FVM) enforces a different property, local conservation, with the same
absence of a penalty. For a control volume $V_i$ with faces $f$, numerical flux $\mathbf{F}_f$, and
outward unit normal $\hat{\mathbf{n}}_f$,
\begin{equation}
  \label{eq:ch37-fvm-cons}
  \frac{\partial}{\partial t} \int_{V_i} u \, dx + \sum_{f \in \partial V_i} \mathbf{F}_f \cdot \hat{\mathbf{n}}_f = 0 .
\end{equation}
Boundary conditions enter by fixing the exterior, or ghost-cell, states from which the boundary-face
flux is computed \index{ghost cell}\index{finite volume method (FVM)!conservative update}. The scope of the guarantee deserves precision, since it is easy
to oversell. Summing \eqref{eq:ch37-fvm-cons} over all cells makes every interior face appear twice
with opposite normals, so the total mass in $\Omega$ changes only through the fluxes the boundary
model itself admits: there is no numerical leakage in the sense of a spurious source, and that holds
to the arithmetic precision of the update regardless of how accurate the flux formula is
\cite{leveque2002finite}\index{conservation!discrete versus learned}. An inconsistent boundary state
yields an exactly conserved wrong answer. The statement again presupposes a mesh with faces on
$\partial\Omega$.

Spectral methods apply the same idea with global bases
\index{spectral method!boundary-conforming basis}: Chebyshev or Legendre polynomials, or spherical
harmonics, whose individual terms can be chosen to satisfy the condition, so that no linear
combination of them can violate it. The Fourier case is the instructive limit, because there the
basis \emph{is} the boundary condition: periodicity on $\mathbb{T}^d$ is built into the trigonometric
system and no trace constraint is left to impose---the assumption behind Part~I's spectral theory
and behind \eqref{eq:fno_layer}. Operational numerical weather prediction (NWP) dynamical
cores such as the European Centre for Medium-Range Weather Forecasts (ECMWF) Integrated Forecasting
System (IFS) truncate a spherical-harmonic expansion horizontally: no lateral boundary, the
surface and the model top closing the vertical column
\index{numerical weather prediction (NWP)!spectral-transform dynamical core}. The cost is
geometry: a global basis conforms to a box, a sphere, or a mapped image of one, and anything else
takes substructuring or penalty machinery.

The lesson that transfers is narrower than it looks. Hard enforcement is free at run time and
expensive per geometry---meshing a new vehicle, generating boundary-layer prisms, re-deriving a
mapping for a new ventricle---and that per-geometry cost, not the linear solve, is what the
multi-query workflows of this part cannot afford
\index{mesh generation!as the recurring cost}. Classical schemes bind exactness to conformity;
operator learning drops the conforming mesh, and with it the source of the exactness.

% -----------------------------------------------------------------------------
\section{Boundary Enforcement in Neural Operators}
\label{sec:ch37-operator}

Neither mechanism below is needed because operators are inaccurate; both are needed because
operators are unconstrained. A generic layer is pointwise lifting and projection around a nonlocal
integral operator, \eqref{eq:no_composition}, and on a uniform grid the spectral specialization
\eqref{eq:fno_layer} makes every one of those operations equivariant to translations on
$\mathbb{T}^d$---the property the multiplier form \eqref{eq:ch05-multiplier} buys---so a network fed
only field values cannot tell a feature at $x$ from the same feature at $x+h$. Nothing in the
composition restricts the trace of the output, whose regularity is not assumed to exceed $L^2(\Omega)$,
where the trace is undefined.

The canonical benchmark makes the exposure concrete. Darcy flow on the unit square,
\eqref{eq:ch25-darcy}, carries homogeneous Dirichlet data while the Fourier neural operator (FNO)
lives naturally on $\mathbb{T}^2$: either the domain is treated as periodic and the mismatch with
$u|_{\partial\Omega}=0$ is absorbed into the learned map, or the field is odd-extended
\index{Darcy flow!boundary condition convention}\index{Fourier neural operator (FNO)!periodic domain and boundary}.
The two define different operators, and a specification silent about which was taken leaves the
dominant error term unnamed. In the DeepONet form \eqref{eq:ch04-deeponet} the same difficulty is one
line of algebra: the trace of $\sum_k b_k(a)t_k(y)$ vanishes for every input only if the trunk
functions cancel on $\partial\Omega$ across a whole family of coefficient vectors, cancellations
gradient descent has no reason to arrange
\index{DeepONet!trunk basis and boundary conditions}.

Absent such structure a trained network learns the boundary condition as a regularity of its labels.
Within the training family that works; outside it the claim is statistical, and three features make it
fragile. The boundary has codimension one: an order-one error confined to a single cell layer of
thickness $h$ on a unit square contributes only order $\sqrt{h}$ to the domain $L^2$ norm, so a
surrogate reported at a few percent relative $L^2$ error can be wrong by an order-one amount exactly
where the objective functional is evaluated
\index{error metric!codimension one of the boundary}. The most variable region of a solution field is
typically the one nearest the boundary, which is why the pointwise normalizers \eqref{eq:ch21-normalize}
find small variance exactly where a condition pins the value. And a change of geometry invalidates the
learned trace wholesale: a wall absent from the training family has no labels on it.

The two available mechanisms differ along one axis: where the constraint lives, in the
parameterization or in the objective.
\index{hard boundary condition!distance envelope}
The first, hard enforcement through a distance envelope, satisfies a homogeneous Dirichlet condition
$u|_{\partial\Omega}=0$ on an arbitrary domain by scaling the output with a continuous distance
function $d(x,\partial\Omega)$ and adding the inhomogeneous part separately,
\begin{equation}
  \label{eq:ch37-hard}
  \mathcal{G}_{\text{hard}}(a)(x) = g(x) + d(x, \partial\Omega) \cdot \mathcal{Q}(v_L(x)),
\end{equation}
where $\mathcal{Q}$ is the pointwise projection of \eqref{eq:no_composition}, applied to the final
latent field $v_L(x)$, so that the envelope multiplies the network's own output, and $g$ is a smooth
extension of the boundary data into $\Omega$. Since
$d(x,\partial\Omega)=0$ for all $x\in\partial\Omega$, the condition holds identically for every
parameter vector $\theta$. That is a one-line proof rather than an experiment, and it is the only
point in this chapter at which an operator architecture attains the classical guarantee. Its
preconditions are classical too: $g$ must exist with the regularity the problem asks for, which
requires the trace to be defined on the output space, i.e.\ $H^{s}(\Omega)$ with $s>1/2$, the
condition \chapref{ch:part1-physics-informed-operators-and-sobolev-reg} records for hard constraints
\index{trace operator}\index{lifting!of inhomogeneous boundary data}.
Three qualifications bound what the exactness is worth.

(i) It is exactness relative to whatever $\partial\Omega$ means in the implementation. The envelope
is evaluated at nodes, so the surface it vanishes on is the discrete boundary of the computing box,
not the body modeled; the residual geometric error is the same object that limits
\eqref{eq:ch37-fem-weak} on a curved boundary, only now it is silent. (ii) The construction
constrains values and nothing else. Where $d$ is differentiable its gradient has unit length and is
normal to $\partial\Omega$, so the normal derivative of the product on the boundary is, up to sign,
$\mathcal{Q}(v_L)$ itself: the flux is left entirely free. That is the correct division of labor for
a Dirichlet problem and no assistance at all for a Neumann, zero-flux, or no-penetration condition.
(iii) There is a regularity cost that
\chapref{ch:part1-discretization-invariance-and-continuum-li} is able to price. A distance function
is Lipschitz on $\overline{\Omega}$ and as smooth as $\partial\Omega$ near it, but loses
differentiability on the interior medial axis, the locus of points with more than one nearest
boundary point \index{signed distance function (SDF)}\index{medial axis!regularity of the distance function}---not an exotic object in the geometries of this part, since the mid-surface of a
ventricular wall, the centerline of a channel, and the interior of a layered aquifer all have one.
Multiplying by a Lipschitz envelope caps the Sobolev regularity of the product; feeding the distance
in as an input channel instead puts a non-periodic function on a periodic grid, jumping across the
wrap wherever the modeled boundary is not a face of the box, with regularity
$H^{1/2-\epsilon}$ for every $\epsilon>0$ and no better---the obstruction that chapter records for
non-periodic coordinate channels, and the reason the discretization rate in \eqref{eq:ch06-rate}
fails to reach $N^{-1}$.

The second mechanism is soft: a variational penalty on the boundary residual,
\begin{equation}
  \label{eq:ch37-soft}
  \mathcal{L}_{\text{boundary}}(\theta) = \lambda_{\text{bc}} \int_{\partial\Omega} |\mathcal{B}(\mathcal{G}_\theta(a)) - g|^2 \, ds,
\end{equation}
where $\mathcal{B}$ is whatever boundary operator the problem states---trace for Dirichlet, normal
flux for Neumann, both for no-slip---and $\lambda_{\text{bc}}$ is a weight. This is the boundary term
of the physics-informed objective \eqref{eq:ch07-bc-loss}, used by the Physics-Informed Neural
Operator (PINO) constructions of \chapref{ch:part1-physics-informed-operators-and-sobolev-reg},
specialized here to the $L^2$ surface norm, so the properties established there transfer: the weight
is dimensional, its usable size depends on the nondimensionalization chosen, and the surface integral
must be realized as a quadrature over a node set on $\partial\Omega$ at every training step
\index{Physics-Informed Neural Operator (PINO)!boundary term}\index{Neumann condition}.
The reported benefits of such training belong to the composite objective and are recorded there~\cite{li2021physics,raissi2019physics};
they are benchmark outcomes, and \eqref{eq:ch37-soft} alone buys no constraint. What the two
mechanisms do buy is priced in Section~\ref{sec:ch37-ledger}.

% -----------------------------------------------------------------------------
\section{Structural Routes: Coordinate Channels, Lifting, and Boundary-Condition-Invariant Parameterizations}
\label{sec:ch37-structural}

Between the two mechanisms sits a family of design choices, differing by which object is modified:
the input representation, the output representation, or the space of admissible outputs itself.

On the input side the cheapest intervention is an extra channel. The grid embedding of
\chapref{ch:part2-the-practitioner-s-paradigm-infinite-dimen} appends normalized coordinates to the
input; a distance-to-boundary channel, a signed distance function, an inside/outside mask, or a
boundary-curvature field are the same device with a physically meaningful scalar in place of $x_1$.
What it buys is a location: the network can condition its response on distance to the wall, so a
boundary layer becomes a function of a physical thickness rather than of a pixel count, and a mask
keeps the model from predicting into solid or void regions at all
\index{coordinate channels!distance to boundary}. What it does not buy is enforcement. The lifting
convolution mixes the channel into the hidden representation, exactly as
\chapref{ch:part2-fourier-neural-operators-fno-from-scratch} describes for the coordinate pair, and
the output trace stays unconstrained: the channel makes the right answer easier to fit and nothing
more.

On the output side, \eqref{eq:ch37-hard} is one instance of a lifting scheme---hand the network the
deviation from a field that already satisfies the condition, then add the field back. For a
grid-based model the envelope is applied after the projection; for DeepONet the same idea folds into
the trunk by replacing $t_k$ with $d\cdot\varphi_k$, so every term of \eqref{eq:ch04-deeponet} has
vanishing trace and exactness holds at every query coordinate for free. The scheme's exposure is the
extension $g$, which is a hypothesis about the boundary data: if the data in use are not well
approximated within the regularity class of the chosen extension, the surrogate inherits an
irreducible error equal to that mismatch, and the extension's smoothness caps the smoothness of the
output, hence the rate in \eqref{eq:ch06-rate}. The scheme has an implementation precondition as
well: it presupposes a distance function that can be evaluated, and on a high-dimensional manifold
evaluating one is the difficulty. Where it bites, practitioners use the penalty
\eqref{eq:ch37-soft} instead.

The third family is the one spectral methods already practice: parameterize the unknown so that
every element of the family satisfies the condition---a stream-function or divergence-free form for
incompressible velocity, a basis of functions vanishing on the boundary, a discretization of the
monodomain operator that admits no flux through the organ surface. In operator learning these remain
largely unbuilt, and their status is the analogue of choosing a conforming Galerkin space: exactness
by construction, and a construction redone whenever the geometry changes. The Fourier basis on the
torus is the single case where the construction is free and uniform---and the geometry it is
indifferent to is fixed at $[0,1]^d$, which is why the periodic benchmark family of
\chapref{ch:part2-canonical-pde-benchmarks-and-reference-wor} is at once cheap and silent about this
chapter's question.

Where geometry genuinely varies, graph-based architectures relocate the problem rather than solve it.
In the geometry-informed neural operator (GINO) the surface mesh \emph{is} the domain of the model:
points are lifted to a latent grid, the spectral work happens there, and the decode returns pressure
and WSS at surface vertices, where the quantities of interest live
\cite{li2023geometry}\index{Geometry-Informed Neural Operator (GINO)!surface as domain}. There is
then no exterior ghost cell to fill and no trace to impose from the network's standpoint---the model
is asked nothing about the field inside the body, and nothing about the velocity on the surface,
whose pressure and traction are all it returns. No-slip survives only as a regularity of the labels
on surface nodes, and nothing in the graph neural operator (GNO) kernel \eqref{eq:ch24-gno-kernel}
obliges a predicted wall traction to be consistent with a wall velocity of zero
\index{no-slip condition!in surface-mesh operators}. The geometry must still be encoded, and
\chapref{ch:part2-graph-neural-operators-gno-and-irregular-g} records the cost: a signed-distance and
coordinate-channel pipeline rebuilt for every new computer-aided design tessellation. The same
substitution appears at internal interfaces. A finite volume scheme puts faces on the caprock, and
conservation across them follows from the telescoping that made \eqref{eq:ch37-fvm-cons} exact; an
operator on a fixed grid sees a fault or layer boundary as a jump in a coefficient channel, and
conservation across it is a learned tendency with no discrete identity behind it
\index{interface condition!impermeable layer}.

% -----------------------------------------------------------------------------
\section{The Ledger: Error Accounting, Proven Claims, and Open Ones}
\label{sec:ch37-ledger}

What the two mechanisms of Section~\ref{sec:ch37-operator} buy can be priced, because the
needed estimate is standard. Let $a$ range over the uniformly elliptic class
$\mathcal{A}_{\text{adm}}$ of Assumption~\ref{ass:ellipticity}, let $v\in H^1(\Omega)$ be any trial
field with trace $\gamma(v)$ and residual $r_v\in H^{-1}(\Omega)$, and let
$u=\mathcal{G}_{\text{true}}(a)$. Lift the trace defect with a bounded right inverse of $\gamma$,
apply the coercivity of \eqref{eq:ch03-coercive} to the remainder, and
\begin{equation}
  \label{eq:ch37-trace-stab}
  \big\|v-\mathcal{G}_{\text{true}}(a)\big\|_{H^{1}(\Omega)}
  \;\le\; C\left(\frac{1}{a_{\min}}\,\|r_v\|_{H^{-1}(\Omega)}
  + \big\|\gamma(v)-g\big\|_{H^{1/2}(\partial\Omega)}\right),
\end{equation}
with $C$ depending on $\Omega$ and on the extension but not on $v$
\cite{evans2010partial,quarteroni2008numerical}\index{elliptic stability estimate}.
The a priori estimate \eqref{eq:ch07-apriori} of
\chapref{ch:part1-physics-informed-operators-and-sobolev-reg} is the case in which the second term
is absent, because there $v$ already has the correct trace. Equation \eqref{eq:ch37-trace-stab} says
more
than that boundary errors add to the total: the boundary enters as a term no amount of residual
minimization removes, measured in $H^{1/2}(\partial\Omega)$, a norm strictly stronger than the
$L^2(\partial\Omega)$ norm \eqref{eq:ch37-soft} shrinks---control in the former would take a
bound on a fractional tangential derivative of the trace defect, which the penalty does not attempt.
The envelope \eqref{eq:ch37-hard} is exactly the device that zeroes the term by construction,
leaving only the question of how small the residual is;
raising $\lambda_{\text{bc}}$ instead trades the boundary term against the data term, with the
recognizable symptom of a small measured boundary residual alongside a field wrong near the
wall. The estimate controls an $H^1$ error while the loss reports a relative $L^2$ one, the second
instance after \eqref{eq:ch31-functional} of a minimized metric and a certified metric that differ.

Table~\ref{tab:ch37-enforcement} arranges the six regimes along the axes that have mattered: what the
constraint is a property of, whether it holds exactly, what it costs, how it fails. The
conclusion is the third column read against the fourth. Exactness is relative to a discretized
object in every row where it appears, the three classical rows no less than the envelope row, and
each classical row buys it with an object tied to the domain---a conforming mesh, a flux stencil
on the faces, a basis for the specific geometry. The envelope row is exact by the same one-line
argument but for values only, on the nodes at which $d$ vanishes; the last two rows are not exact
at all. What the three operator rows share is nothing spent per query.

\begin{table}[htbp]
\centering
\small
\caption{Where a boundary condition can live: six enforcement regimes compared.}
\label{tab:ch37-enforcement}
\begin{tabularx}{\textwidth}{@{}>{\raggedright\arraybackslash}p{0.14\textwidth}%
  >{\raggedright\arraybackslash}p{0.16\textwidth}%
  >{\raggedright\arraybackslash}p{0.13\textwidth}%
  >{\raggedright\arraybackslash}p{0.16\textwidth}%
  >{\raggedright\arraybackslash}X@{}}
\toprule
\textbf{Regime} & \textbf{Condition lives in} & \textbf{Exact?} & \textbf{Marginal cost} & \textbf{Dominant failure mode} \\
\midrule
Conforming FEM & trial and test spaces & on the discrete space & mesh per geometry & curved-boundary geometry error \\
FVM with ghost cells & boundary-face flux & conserved averages & ghost-layer construction & non-conforming interfaces \\
Spectral & basis functions & for the modeled condition & basis tied to domain & nonsmooth or moving boundaries \\
Operator plus envelope & output parameterization & on the sampled $\partial\Omega$ & distance evaluation at setup & value-type conditions only \\
Operator plus penalty & objective function & no & boundary quadrature per step & norm mismatch and weight tuning \\
Operator plus labels & training distribution & no & nothing at inference & new geometry or new loading \\
\bottomrule
\end{tabularx}
\end{table}

Sort the remaining claims by status. Proven: the exactness of
\eqref{eq:ch37-hard} with respect to the set on which $d$ is evaluated; the estimate
\eqref{eq:ch37-trace-stab} under Assumption~\ref{ass:ellipticity} with boundary regularity on
$\Omega$; and the disappearance of the question on $\mathbb{T}^d$ and $\mathbb{S}^2$, where
\eqref{eq:ch07-bc-loss} has no surface to integrate over. Observed: that adding
\eqref{eq:ch37-soft} lowers the measured boundary residual, that composite physics-informed training
reduces label requirements and improves behavior above the training resolution~\cite{li2021physics},
and that geometry channels let one trained surface-mesh operator serve the nearby shapes of
\chapref{ch:part3-external-aerodynamics-and-surface-flows}. Not established
anywhere in this book: a theorem that a trained operator satisfies a boundary condition
to a stated tolerance, on a stated class of domains, in a norm where the trace is defined
\index{universal approximation!boundary conditions outside the norm}.
That gap is structural rather than merely unproved. The universality
bound \eqref{eq:ch04-unif-bound} is an $L^2(\Omega)$ estimate, and the trace operator is not a
continuous functional of $L^2(\Omega)$, so nothing prevents a sequence Cauchy in that norm from
violating the condition. Approximation theory and boundary analysis live in different
topologies, which is why exactness in operator learning must be engineered into the architecture, as
in \eqref{eq:ch37-hard}, rather than inherited from a density result.

Two consequences follow, both constraints on practice rather than on style. The first is
diagnostic: a surrogate whose objective lives on the boundary should be scored by the boundary
defect itself, in a norm that can be named, on held-out geometries and loadings. A relative
$L^2$ number on the training family is not evidence about the trace, and the near-zero-variance
locations of \eqref{eq:ch21-normalize} should be checked rather than tolerated, since they mark where
the condition is imposed. The second is institutional. An auditor of a classical computation
inspects a constraint space: the boundary condition is an enumerable property of the discretization,
exactly as the flux stencil of \eqref{eq:ch37-fvm-cons} is auditable. A
learned trace is a statistical claim about a label distribution, and the labels nearest the boundary
are the costliest to generate, because boundary-layer resolution sets the size of the whole mesh. The
region that matters most is where the data are thinnest---the regime the penalty
\eqref{eq:ch37-soft} was written for, and the one in which it is least verifiable
\index{certification!auditable constraint spaces}\index{data availability!boundary labels}.

% -----------------------------------------------------------------------------

Boundary conditions are the clearest instance of a pattern this part has met five times under five
names: a property that a conforming discretization carries by construction becomes, in a learned
operator, either a term to be traded off in an objective or a regularity of the data to be tested
for. The pattern is not confined to the trace. Conservation of mass, momentum, and energy fails in
the same silent way, and so does reliability outside the training distribution, where a classical
solver diverges, violates a time-step restriction, or refuses to converge while a network returns a
smooth, plausible, wrong field. \chapref{ch:part3-synthesis-matrix-and-real-world-deployment}
assembles these into a single cross-domain ledger---the five application chapters with the cost
ratio each row can form and the two rows that cannot, their fidelity metrics, and the rigorous
statement Part~I supplies---and then takes up five roadblocks between a published surrogate and a
deployed one: error estimation without a reference solver, exact conservation, the price of labels,
distribution shift, and clinical clearance. This chapter supplies the template for that synthesis: for each guarantee, ask
which object carries it, whether it is exact, what it costs, and what would have to be measured to
believe it.
